\documentclass[../BTL.tex]{subfiles}
\begin{document}


\section{Mục tiêu và phạm vi đề tài}
\label{section:1.2}

Mục tiêu của đề tài là nghiên cứu, cài đặt từ gốc (from scratch) và so sánh các mô hình Deep Learning cho bài toán phân loại bình luận độc hại, nhằm tìm hiểu sâu về kiến trúc và hành vi của từng mô hình. Việc cài đặt từ đầu giúp hiểu rõ cơ chế hoạt động bên trong thay vì chỉ sử dụng các thư viện có sẵn.

Về phạm vi nghiên cứu, đề tài tập trung vào việc cài đặt và so sánh ba mô hình Deep Learning chính, bao gồm: Recurrent Convolutional Neural Network (RCNN) với phiên bản cải tiến EnhancedRCNN~\cite{lai2015rcnn}, Attention BiLSTM~\cite{yang2016hierarchical}, và Transformer Encoder~\cite{vaswani2017transformer}. Tất cả các mô hình này đều được cài đặt từ đầu mà không sử dụng các thư viện có sẵn như nn.LSTM, nn.GRU, hay nn.Transformer của PyTorch, nhằm mục đích thể hiện sự hiểu biết sâu về kiến trúc bên trong của các mô hình. Việc cài đặt từ gốc bao gồm xây dựng LSTM Cell và GRU Cell từ các phép toán ma trận cơ bản, cơ chế Attention từ đầu, các lớp Self-Attention và Multi-Head Attention, cũng như các chiến lược Pooling có trọng số (masked mean/max pooling).

Đề tài cũng áp dụng nhiều kỹ thuật tối ưu hóa: Masked Pooling, Focal Loss~\cite{lin2017focal}, Label Smoothing~\cite{szegedy2016rethinking}, Gradient Accumulation, AdamW Optimizer~\cite{loshchilov2017adamw}, Gradient Clipping~\cite{graves2012supervised}, OneCycleLR Scheduler~\cite{smith2018superconvergence}, Early Stopping dựa trên ROC-AUC, và kỹ thuật tìm kiếm Optimal Thresholds.

Để thực hiện các thử nghiệm, đề tài sử dụng hai bộ dữ liệu: Jigsaw Toxic Comment Classification Challenge~\cite{kaggleJigsaw} và HateXplain~\cite{hatexplain2021}. Chi tiết về các bộ dữ liệu được trình bày trong Chương 3.

\section{Bố cục báo cáo}
\label{section:1.4}
Báo cáo được tổ chức thành bốn chương. Chương 1 giới thiệu tổng quan về mục tiêu và phạm vi đề tài. Chương 2 trình bày cơ sở lý thuyết về các mô hình và kỹ thuật tối ưu hóa được sử dụng. Chương 3 mô tả các thử nghiệm, kết quả đánh giá và phân tích. Chương 4 tổng kết kết quả đạt được và đề xuất hướng phát triển.

\end{document}
